\section{Rough kernels and the annular endpoint errors}
\label{sec:rough}

The comparison \eqref{eq:hard-smooth} leaves one error kernel at each
annular endpoint.  Estimating these kernels separately would cost the
number of endpoints.  We instead collect endpoints at comparable scales.
Disjointness of the annuli gives a Gaussian bound for each resulting
kernel, independent of the number of its discontinuities.  The frequency
estimate from the preceding section then controls the groups with few
discontinuities; the remaining groups are few enough to estimate
pointwise.

\subsection{A pointwise Gaussian estimate}

For a locally integrable function on $\R^2$, let $M_x$ and $M_y$ denote
the centered Hardy--Littlewood maximal operators in the first and second
coordinates.  For $F,G\in L^3(\R^2)$, put
\begin{equation}
 \mathcal M_{F,G}(x,y)
 =\bigl(M_x(|F|^{5/2})(x,y)\bigr)^{2/5}
  \bigl(M_y(|G|^{5/2})(x,y)\bigr)^{2/5}.
 \label{eq:rough-maximal-envelope}
\end{equation}
The one-dimensional maximal inequality on $L^{6/5}$, followed by
H\"older's inequality on $\R^2$, gives
\begin{equation}
 \|\mathcal M_{F,G}\|_{3/2}
 \le C\|F\|_3\|G\|_3.
 \label{eq:rough-maximal-norm}
\end{equation}

\begin{lemma}[Gaussian envelope]
\label{lem:gaussian-envelope}
Suppose that $c>0$ and that a measurable kernel $f$ satisfies
$|f(z)|\le e^{-cz^2}w(z)$, where $w\ge0$ and $\|w\|_5\le\delta$.
Then, for every $L>0$,
\begin{equation}
 \int_{\R}|F(x+t,y)G(x,y+t)|\,|\Dil_L f(t)|\,\dd t
 \le C_c\delta\,\mathcal M_{F,G}(x,y)
 \label{eq:gaussian-envelope}
\end{equation}
at every point where the right side is finite.  In particular, a bound
$|f(z)|\le A e^{-cz^2}$ gives \eqref{eq:gaussian-envelope} with
$C_c\delta$ replaced by $C_cA$.
\end{lemma}

\begin{proof}
After setting $t=Lz$, apply H\"older's inequality with exponents
$5/2,5/2,5$ to
\[
 |F(x+Lz,y)|e^{-cz^2/2},\qquad
 |G(x,y+Lz)|e^{-cz^2/2},\qquad w(z).
\]
The first resulting integral is
$\int |F(x+Lz,y)|^{5/2}e^{-5cz^2/4}\,\dd z$.
Splitting $\R$ into $|z|\le1$ and
$2^j<|z|\le2^{j+1}$ bounds it by
$C_cM_x(|F|^{5/2})(x,y)$, uniformly in $L$.
The second integral has the corresponding bound with $G$ and $M_y$.
This proves the first assertion.  For the last assertion, retain half
the Gaussian decay in $w$, whose $L^5$ norm is then at most $C_cA$.
\end{proof}

We will use this estimate inside the triangular form.  Recall that
$G_1(b,c)=F(b+c,-b)$ and $G_2(c,a)=G(-a,c+a)$, and that these changes
of coordinates preserve the $L^3$ norms.  With $t=a+b+c$, the absolute
value of the $c$-integral is bounded by
$C_c\delta\mathcal M_{F,G}(-a,-b)$.
Consequently, integrating against $|G_0(a,b)|$ costs at most
$C_c\delta\prod_{v=0}^2\|G_v\|_3$.
The kernel, its dilation, and its envelope may depend on $(a,b)$:
the same assertion holds whenever the bound for $\delta$ is uniform.

\subsection{Families with a controlled number of jumps}

The next proposition turns the frequency estimate into an estimate for
piecewise smooth kernels.  The pointwise bound on the number of jumps
at one scale and the bound on their total number play different roles:
they will give, respectively, the uniform bound and the square-sum
bound for the frequency coefficients.

\begin{proposition}[Rough kernel families]
\label{prop:rough-family}
Fix $A\ge1$ and $c>0$.  Let $n\ge1$, and let $\mathcal I\subset\mathbb Z$
be finite.  For $k\in\mathcal I$, let $K_{a,b,k}\colon\R\to\C$ and
$m_{a,b,k}\in\{0,1,2,\ldots\}$ be finite-valued step choices in
$(a,b)$, with $K_{a,b,k}=0$ whenever $m_{a,b,k}=0$.
When $m_{a,b,k}\ge1$, suppose that $K_{a,b,k}$ is $C^1$ outside a set
of at most $A m_{a,b,k}$ points and that, outside that set,
\begin{equation}
 |K_{a,b,k}(z)|+|K_{a,b,k}'(z)|\le A e^{-cz^2}.
 \label{eq:rough-family-decay}
\end{equation}
Assume also that
\begin{equation}
 \int_{\R}z^\ell K_{a,b,k}(z)\,\dd z=0
 \quad(\ell=0,1,2),
 \qquad
 \sum_{k\in\mathcal I}m_{a,b,k}\le An,
 \quad m_{a,b,k}\le A\sqrt n.
 \label{eq:rough-family-hypotheses}
\end{equation}
For
\[
 \mathcal K_{a,b}(t)
 =\sum_{k\in\mathcal I}\Dil_{2^k}K_{a,b,k}(t),
\]
and smooth compactly supported $G_0,G_1,G_2$, one has
\begin{equation}
 |\Lambda_{\mathcal K}(G_0,G_1,G_2)|
 \le C_{A,c}\sqrt n\log(2+n)\prod_{v=0}^2\|G_v\|_3.
 \label{eq:rough-family-bound}
\end{equation}
\end{proposition}

\begin{proof}
We first obtain Fourier estimates for a single kernel $K=K_{a,b,k}$;
write $m=m_{a,b,k}$.  All constants below depend only on $A,c$.
There is nothing to prove for $m=0$, so suppose that $m\ge1$.

\paragraph{A Gaussian-weighted primitive.}
We seek a representation $K=(g_{3q}f)'''$ for a suitable fixed $q$:
the Fourier modes of $f$ then produce exactly the differentiated
Gaussian windows controlled by Corollary~\ref{cor:frequency-continuum}.
Define
\[
 U(z)=\frac12\int_{-\infty}^z(z-u)^2K(u)\,\dd u.
\]
The three moment conditions imply, for $j=0,1,2$,
\[
 U^{(j)}(z)
 =-\frac1{(2-j)!}\int_z^\infty(z-u)^{2-j}K(u)\,\dd u.
\]
Use the defining integral when $z\le0$ and this tail integral when
$z\ge0$.  Gaussian tail estimates give
$|U^{(j)}(z)|\le C e^{-c_1z^2}$ for $j=0,1,2$, with $c_1>0$ fixed.
Moreover, $U'''=K$ as a weak derivative.
Choose $q>0$ sufficiently large, depending only on $A,c$, and put
\[
 f(z)=\frac{U(z)}{g_{3q}(z)}.
\]
The quotient and product rules, with a slight reduction in the
Gaussian exponent, give
\begin{equation}
 |f^{(j)}(z)|\le C e^{-c_2z^2}\quad(0\le j\le3).
 \label{eq:rough-f-decay}
\end{equation}
Here and below derivatives of order three are weak derivatives, or
their almost-everywhere representatives.  The functions $f,f',f''$
are continuous.  Between the exceptional points, the same calculation
using $K'$ bounds the classical derivative of $f'''$ by
$Ce^{-c_2z^2}$.  Its jumps are bounded by $C$, and there are at most
$Am$ of them.  Thus, for $0\le j\le3$,
\begin{equation}
 \|f^{(j)}\|_\infty\le C,
 \qquad \operatorname{Var}(f^{(j)})\le Cm.
 \label{eq:rough-f-bv}
\end{equation}
The same bounds hold with zero right sides when $m=0$.

\paragraph{Fourier bounds.}
Use the normalization
$\widehat f(\xi)=(2\pi)^{-1}\int f(z)e^{-i\xi z}\,\dd z$.
Let
\[
 \Xi_1=\{\xi:|\xi|<2\},\qquad
 \Xi_D=\{\xi:D\le|\xi|<2D\}
 \quad(D=2,4,8,\ldots).
\]
We claim that
\begin{equation}
 \sup_{\xi\in\Xi_D}|D^3\widehat f(\xi)|\le\frac{Cm}{D},
 \qquad
 \int_{\Xi_D}|D^3\widehat f(\xi)|^2\,\dd\xi
 \le\frac{Cm}{D}.
 \label{eq:rough-fourier-bounds}
\end{equation}
For $D\ge2$, the distributional derivative of $f'''$ is a finite
measure of total variation at most $Cm$.  Taking its Fourier transform
gives $|\xi|^4|\widehat f(\xi)|\le Cm$, proving the first bound.
For the second, the translation inequality for functions of bounded
variation gives
\[
 \|f'''(\,\cdot+D^{-1})-f'''\|_1\le\frac{Cm}{D}.
\]
Indeed, integrate the absolute derivative measure over each traversed
interval and apply Fubini's theorem.  By the uniform supremum bound
in \eqref{eq:rough-f-bv}, the squared $L^2$ norm of this difference is
also at most $Cm/D$.  Plancherel's identity and
$\widehat{f'''}(\xi)=(i\xi)^3\widehat f(\xi)$ now yield the second
bound, since $|e^{i\xi/D}-1|\ge2\sin(1/2)$ on $\Xi_D$.
When $D=1$, both assertions follow from
\eqref{eq:rough-f-decay} and $m\ge1$.

\paragraph{A finite frequency cutoff.}
Choose $\chi\in C_c^\infty(\R)$ supported on $[-2,2]$ and equal to one
on $[-1,1]$, and set
\[
 P_*=(2+n)^8,\qquad
 \widehat{f_*}(\xi)=\widehat f(\xi)\chi(\xi/P_*),\qquad
 K_*=(g_{3q}f_*)'''.
\]
We show that this replacement has a summable pointwise error.
Let $\eta$ be the Schwartz convolution kernel of the multiplier
$\chi$, normalized so that $\int\eta=1$, and let
$\eta_{P_*}(z)=P_*\eta(P_*z)$.  For $0\le j\le3$,
$f_*^{(j)}=f^{(j)}*\eta_{P_*}$.  The translation inequality gives
\[
 \|f^{(j)}-f_*^{(j)}\|_1
 \le\operatorname{Var}(f^{(j)})
      \int |u|\,|\eta_{P_*}(u)|\,\dd u
 \le\frac{Cm}{P_*}.
\]
Their supremum norms are uniformly bounded, by
\eqref{eq:rough-f-bv} and the $L^1$ norm of $\eta$.  Interpolation
therefore gives
\[
 \|f^{(j)}-f_*^{(j)}\|_5
 \le C(m/P_*)^{1/5}.
\]
Since every derivative of $g_{3q}$ is bounded by a constant times a
slightly wider Gaussian, the product rule shows that
\begin{equation}
 |K(z)-K_*(z)|\le e^{-c_3z^2}w(z),
 \qquad \|w\|_5\le C(m/P_*)^{1/5},
 \label{eq:rough-cutoff-error}
\end{equation}
for a nonnegative $w$.  This calculation also holds for weak
derivatives, so no additional terms occur at the jumps of $K$.

Apply this construction separately to every $K_{a,b,k}$.
By \eqref{eq:rough-family-hypotheses},
$(m_{a,b,k}/P_*)^{1/5}\le C/n$ whenever $m_{a,b,k}\ne0$, and at
each $(a,b)$ there are at most $An$ such indices.  Lemma
\ref{lem:gaussian-envelope} bounds the sum of all replacement errors
pointwise by $C\mathcal M_{F,G}(-a,-b)$ after integration in $c$.
Their contribution to $\Lambda_{\mathcal K}$ is consequently at most
$C\prod_v\|G_v\|_3$.

\paragraph{Applying the frequency estimate.}
It remains to estimate the kernels $K_*$.  For each dyadic $D$ use
the kernels from Corollary~\ref{cor:frequency-continuum}, namely
\[
 \ell_\xi(z)=D^{-1}\frac{\dd}{\dd z}
                 \bigl(g_q(z)e^{i\xi z}\bigr),
 \qquad H_\xi=\ell_\xi*\ell_\xi*\ell_\xi.
\]
Gaussian convolution and modulation give
\[
 (g_{3q}(z)e^{i\xi z})'''=D^3H_\xi(z).
\]
For the block $\Xi_D$, the contribution to $K_{a,b,k,*}$ is therefore
\[
 \int_{\Xi_D}c_{a,b,k}(\xi)H_\xi(z)\,\dd\xi,
 \qquad
 c_{a,b,k}(\xi)
 =D^3\widehat f_{a,b,k}(\xi)\chi(\xi/P_*).
\]
Set $\mu_{a,b,k}(\xi)=D c_{a,b,k}(\xi)/\sqrt n$.
The two estimates in \eqref{eq:rough-fourier-bounds} give precisely
\[
 |\mu_{a,b,k}(\xi)|
 \le\frac{Cm_{a,b,k}}{\sqrt n}\le C,
 \qquad
 \sum_k\int_{\Xi_D}|\mu_{a,b,k}(\xi)|^2\frac{\dd\xi}{D}
 \le C\sum_k\frac{m_{a,b,k}}n\le C.
\]
All transformations were applied separately to finitely many kernel
values.  The coefficients $\mu_{a,b,k}$ therefore remain step choices
on a common finite rectangular refinement of the original partitions.
The block kernel is $\sqrt n\int_{\Xi_D}\mu_{a,b,k}(\xi)
H_\xi\,\dd\xi/D$.  Thus Corollary~\ref{cor:frequency-continuum},
with measure $\dd\xi/D$, bounds this block by
$C\sqrt n\prod_v\|G_v\|_3$.
Only $O(\log(2+n))$ blocks meet the support of the cutoff.
Summing their bounds and adding the replacement error proves
\eqref{eq:rough-family-bound}.
\end{proof}

\subsection{Removing the first moment of an odd kernel}

The endpoint kernels are odd but need not have vanishing first moment.
A short dilation identity reduces them to the preceding proposition.

\begin{corollary}[Odd kernel families]
\label{cor:odd-family}
Proposition~\ref{prop:rough-family} remains valid if the three moment
conditions in \eqref{eq:rough-family-hypotheses} are replaced by the
assumption that every $K_{a,b,k}$ is odd almost everywhere.
\end{corollary}

\begin{proof}
Choose a real, smooth, compactly supported odd function $\varphi$ with
$\int z\varphi(z)\,\dd z=1$.  Put
\[
 u_{a,b,k}=\int zK_{a,b,k}(z)\,\dd z.
\]
These coefficients are uniformly bounded and vanish when
$m_{a,b,k}=0$.  The kernels $K_{a,b,k}-u_{a,b,k}\varphi$ satisfy all
three moment conditions and the other hypotheses of
Proposition~\ref{prop:rough-family}, with enlarged fixed constants
and unchanged counts.  It remains to bound the subtracted terms.

Define $\Psi=\varphi-2^{-1}\Dil_2\varphi$.
Oddness gives its zeroth and second moments, and its first moment is
$1-2^{-1}\cdot2=0$.  For every integer $J\ge1$,
\[
 \Dil_{2^k}\varphi
 =\sum_{j=0}^{J-1}2^{-j}\Dil_{2^{k+j}}\Psi
   +2^{-J}\Dil_{2^{k+J}}\varphi.
\]
For each fixed $j$, apply Proposition~\ref{prop:rough-family} to the
shifted scales $k+j$ with kernels $u_{a,b,k}\Psi$ and the same shifted
counts $m_{a,b,k}$.  The constants are independent of $j$, and the
outer coefficients $2^{-j}$ are summable.  The remainder contributes
at most $C2^{-J}n\prod_v\|G_v\|_3$ by Lemma
\ref{lem:gaussian-envelope}, since there are at most $An$ nonzero
counts at any output point.  Letting $J\to\infty$ proves the claim.
\end{proof}

\subsection{Grouping the endpoint errors}

We now verify that the errors in \eqref{eq:hard-smooth} have the
structure just proved.  Recall
\[
 E_\rho(t)
 =\frac{P(|t|/\rho)-c_0\mathbf1_{\{|t|>\rho\}}}{t},
 \qquad
 P(v)=2\int_0^v\psi(u)\,\dd u,
 \qquad \psi=-g_3'''.
\]
We take the continuous extension at $t=0$; values at the jump points
do not affect any integral.

\begin{proposition}[Endpoint errors]
\label{prop:endpoint-errors}
At each $(a,b)$, choose at most $n$ pairwise disjoint open radius
intervals $(\varepsilon_j,R_j)$ with $0<\varepsilon_j<R_j<\infty$,
and complex coefficients
$\beta_j$ with $|\beta_j|\le1$.  Suppose that these choices form a
finite-valued step function of $(a,b)$.  Then the kernel
\[
 \mathcal E_{a,b}(t)
 =\sum_j\beta_j
       \bigl(E_{\varepsilon_j}(t)-E_{R_j}(t)\bigr)
\]
satisfies
\begin{equation}
 |\Lambda_{\mathcal E}(G_0,G_1,G_2)|
 \le C\sqrt n\log(2+n)\prod_{v=0}^2\|G_v\|_3
 \label{eq:endpoint-error-bound}
\end{equation}
for smooth compactly supported inputs.
\end{proposition}

\begin{proof}
Place each endpoint $\rho$ in the unique half-open dyadic interval
$[L_k,2L_k)$, where $L_k=2^k$.  Let $m_{a,b,k}$ count these endpoints,
with multiplicity, and collect their contributions into
$\Dil_{L_k}K_{a,b,k}$.  Thus $K_{a,b,k}$ is a linear combination of
\begin{equation}
 e_\alpha(z)
 =\frac{P(|z|/\alpha)-c_0\mathbf1_{\{|z|>\alpha\}}}{z},
 \qquad 1\le\alpha<2,
 \label{eq:dimensionless-endpoint}
\end{equation}
with coefficient $\beta_j$ at a lower endpoint and $-\beta_j$ at its
upper endpoint.  Every kernel is odd, its possible jump points are
the $\pm\alpha$, and
\begin{equation}
 \sum_km_{a,b,k}\le2n.
 \label{eq:endpoint-count}
\end{equation}

The important assertion is the uniform bound
\begin{equation}
 |K_{a,b,k}(z)|+|K_{a,b,k}'(z)|\le Ce^{-cz^2}
 \label{eq:paired-endpoint-bound}
\end{equation}
away from at most $2m_{a,b,k}$ points, with constants independent of
$m_{a,b,k}$.  Fix $(a,b)$ and $k$.  Pair the two endpoints of every
annulus whose endpoints both lie in $[L_k,2L_k)$.  After scaling, its
contribution is $\beta_j(e_{\alpha_j}-e_{\gamma_j})$ with
$1\le\alpha_j<\gamma_j<2$.
The intervals $(\alpha_j,\gamma_j)$ are disjoint, so
\begin{equation}
 \sum_j(\gamma_j-\alpha_j)\le1.
 \label{eq:paired-length}
\end{equation}
There are at most two unpaired endpoint contributions.  Their
intervals must satisfy, respectively,
\[
 \varepsilon_j<L_k\le R_j<2L_k,
 \qquad\text{or}\qquad
 L_k\le\varepsilon_j<2L_k\le R_j.
\]
Disjointness allows at most one interval of each type.  An interval
with $\varepsilon_j<L_k$ and $R_j\ge2L_k$ has no endpoint in this
group.  Figure~\ref{fig:endpoint-pairing} illustrates the decomposition.

\begin{figure}[t]
\centering
\begin{tikzpicture}[x=1cm,y=1cm,>=Stealth]
 \fill[gray!10] (1,-.25) rectangle (8,1.55);
 \draw[->] (0,0)--(9,0) node[right] {$\rho$};
 \draw[densely dashed] (1,-.25)--(1,.93);
 \draw[densely dashed] (8,-.25)--(8,.93);
 \node[below] at (1,0) {$L_k$};
 \node[below] at (8,0) {$2L_k$};
 \draw[thick] (.25,.55)--(1.65,.55);
 \draw[thick] (2.3,.55)--(3.55,.55);
 \draw[thick] (4.15,.55)--(4.8,.55);
 \draw[thick] (5.45,.55)--(6.4,.55);
 \draw[thick] (7.15,.55)--(8.75,.55);
 \foreach \x in {.25,1.65,2.3,3.55,4.15,4.8,5.45,6.4,7.15,8.75}
   \fill (\x,.55) circle (1.5pt);
 \draw[decorate,decoration={brace,amplitude=4pt}]
   (2.3,.88)--(6.4,.88);
 \node[above] at (4.35,1.04) {paired endpoints};
 \node[above,align=center,font=\small] at (1.05,1.03) {one unpaired};
 \node[above,align=center,font=\small] at (7.95,1.03) {one unpaired};
\end{tikzpicture}
\caption{Endpoint grouping inside one dyadic interval.  Intervals with
both endpoints in the group contribute differences; their total
length is at most $L_k$.  Only intervals reaching a boundary of the
dyadic interval can leave unpaired endpoints.  The picture is schematic.}
\label{fig:endpoint-pairing}
\end{figure}

For $|z|\le3$, first consider the smooth term
$s_\alpha(z)=P(|z|/\alpha)/z$ in \eqref{eq:dimensionless-endpoint}.
The function $P$ has a smooth even extension and vanishes quadratically
at zero.  Hence $s_\alpha$ extends smoothly across $z=0$, and
\[
 \sup_{\substack{|z|\le3\\1\le\alpha\le2}}
 \bigl(|\partial_\alpha s_\alpha(z)|
       +|\partial_z\partial_\alpha s_\alpha(z)|\bigr)<\infty.
\]
The fundamental theorem of calculus in $\alpha$, followed by
\eqref{eq:paired-length}, bounds the sum of the paired smooth terms
and their $z$-derivatives independently of their number.
For the hard terms, the paired difference is
\[
 -\frac{c_0}{z}
    \bigl(\mathbf1_{\{|z|>\alpha_j\}}
          -\mathbf1_{\{|z|>\gamma_j\}}\bigr).
\]
Away from endpoints, at most one such difference is nonzero at any
fixed $|z|$, because the intervals $(\alpha_j,\gamma_j)$ are disjoint.
On its support $|z|\ge1$, so its size and its classical derivative
are uniformly bounded.  The at most two unpaired terms have the
same bounds on this compact range.

For $|z|>3$, all indicators in \eqref{eq:dimensionless-endpoint}
equal one, and
\[
 P(|z|/\alpha)-c_0=-2g_3''(|z|/\alpha).
\]
Consequently $e_\alpha$, $\partial_z e_\alpha$,
$\partial_\alpha e_\alpha$, and
$\partial_z\partial_\alpha e_\alpha$ are bounded by
$Ce^{-cz^2}$, uniformly for $1\le\alpha\le2$.
The paired-length estimate and the two unpaired contributions prove
the same bound for their sum.  Combining the compact and tail
estimates proves \eqref{eq:paired-endpoint-bound}.

We finish by separating the groups according to their counts.
At each output point, \eqref{eq:endpoint-count} implies that there
are at most $2\sqrt n$ indices with $m_{a,b,k}>\sqrt n$.
The uniform Gaussian bound \eqref{eq:paired-endpoint-bound} and
Lemma~\ref{lem:gaussian-envelope} bound all these groups together,
after integration in $c$, by
$C\sqrt n\mathcal M_{F,G}(-a,-b)$.
Their trilinear contribution is at most
$C\sqrt n\prod_v\|G_v\|_3$.

Set the kernels and counts of those groups to zero.  The remaining
family satisfies \eqref{eq:rough-family-decay}, has odd kernels,
and has counts satisfying
$\sum_km_{a,b,k}\le2n$ and $m_{a,b,k}\le\sqrt n$.
Corollary~\ref{cor:odd-family} applies and proves
\eqref{eq:endpoint-error-bound}.  Every grouping and deletion is made
at the output point, so the proof imposes no common choice of
endpoints or dyadic groups at different output points.
\end{proof}
